%% Cover letter for Applied Stochastic Models in Business and Industry
%% Compiled with pdflatex
\documentclass[10pt,a4paper]{article}
\usepackage[top=20mm,bottom=18mm,left=24mm,right=24mm]{geometry}
\usepackage{newtxtext}
\usepackage[hidelinks]{hyperref}
\usepackage{parskip}
\pagestyle{empty}
\setlength{\parskip}{5pt}

\begin{document}

\begin{flushright}
Junjie Zhang\\
Shanghai Academy of Global Governance and Area Studies\\
School of Economics and Finance\\
Shanghai International Studies University, Shanghai 201620, China\\
junjiezhang2024@shisu.edu.cn \quad ORCID 0009-0004-8821-4018
\end{flushright}

8 October 2026

To the Editor in Chief, \emph{Applied Stochastic Models in Business and Industry}

\textbf{Re: submission of a Research Article}\\
\textbf{``From Joint Probability Forecasts to Selective Review: Decision Value and Calibration Sensitivity in Electricity Monitoring''}

Dear Editor,

Please consider the enclosed manuscript for publication as a Research Article in \emph{Applied Stochastic Models in Business and Industry}.

\textbf{The problem.} A monitoring desk that receives hourly joint risk forecasts can hold a well calibrated probability and still review the wrong part of a developing episode. Review capacity is limited, so the number of delivery windows that can be examined is small, and the first hours of a persistent process carry more decision value than its later hours. The manuscript treats this as a stochastic decision problem rather than a forecasting exercise: it asks how probability diagnostics translate into selections, what those selections cost in workload, and how far small changes in marginal calibration move the resulting review list. The application is synchronous high net load in Germany and Luxembourg, France and Belgium over 17,400 scored hours and 71 processes, with Danish replacement-zone analyses delimiting the transfer evidence.

\textbf{What the study adds.} Three features distinguish it from the forecasting literature it builds on. First, the evaluation separates event hour hits, selected delivery windows and coverage of a process's first six hours, so that discovery and repetition are scored as different outcomes. Second, a controlled marginal-map and dependence-family comparison is followed by probability perturbations that are propagated all the way through ranking and stateful allocation, which shows that selection stability, rather than average score, is the binding constraint: local marginal stress changes 16.2 percent of ordinary and 28.0 percent of stateful baseline selections. Third, a forecast built for the early-event target, evaluated against an otherwise identical hourly-target model with the same covariates, quantifies the value of aligning the statistical target with the review objective: 43 processes covered against 38 and 35. Selective gating moves selections from 1,450 to 618 while keeping 322 of the 341 event hours that daily Top2 finds, and early coverage moves from 35 to 32 processes, so the workload saving and its discovery cost are reported together rather than traded off silently.

\textbf{Why this journal.} The work is a stochastic-model evaluation with an explicit industrial decision: probability calibration, dependence structure, threshold sensitivity and a cost-based selection contract are analysed together in an operational monitoring setting, which is the combination the journal publishes. It extends the applied literature on multivariate electricity and energy forecasting that the journal has recently carried, and it keeps the stochastic argument central: a coupling bound for fixed copulas, sufficient margins for ordinary selection stability, and known-mechanism simulations that separate gate crossings from capacity competition.

\textbf{Statements.} The research code and the derived evaluation data are openly available in Zenodo, version 1.1.2, DOI \href{https://doi.org/10.5281/zenodo.23230336}{10.5281/zenodo.23230336}. The manuscript is original and is not under consideration elsewhere. No external funding supported this research and I declare no competing interests. The study uses public records and simulations, with no human participants or personal data. Supporting Information provides full specifications, result tables, Figures S1 and S2, and Tables S1 to S12.

Thank you for considering this submission; I would be pleased to respond to any questions.

Yours sincerely,

\vspace{8pt}
Junjie Zhang\\
Shanghai International Studies University\\
junjiezhang2024@shisu.edu.cn

\end{document}
